\section{Conclusions and Future Work}
The evaluation establishes that a single Multi-Agent Reinforcement Learning (MARL) policy can generalize to previously unseen network backbones and enhance the deployed baseline at the packet level. This policy is trained using an analytical model and evaluated at the packet level. The research shows that multi-agent reinforcement learning combined with a graph neural network (GNN) routing protocol outperforms existing routing protocols such as OSPF and single-agent policies. In overload scenarios, the policy reduced packet loss from 14.1\% to 3.4\% on GÉANT and from 27.3\% to 10.6\% on Germany50 compared to OSPF. Furthermore, the results show that MARL with GNN achieves superior performance on medium to large networks, such as GEANT and Germany50. However, on the smallest topology, it does not reduce packet loss below OSPF. In the feasible range, the extra headroom comes at the cost of a slight increase in packet loss.

For future work, we recommend training the model at larger scales with packet-level data and evaluating more network topologies and traffic patterns. We also suggest extending the experiments and increasing compute power to achieve comprehensive results. In addition, this study used the actor-critic Proximal Policy Optimization algorithm for rewards, so other algorithms could be explored in the future. Testing the model on real hardware routers could also help move toward a production-ready model.
